% Unit 129 terjemahan Bahasa Indonesia dari chapter9.tex baris 1369--1456.
% Sumber: Methods of Algebra, Volume 2, commit
% 9a5803ff77dd3257484cb177f851a73770a59dd3 (CC BY 4.0).
% stable-unit-id: o014.aljabr2.chapter9.reconstruction-b
% Cakupan: sisa Bagian 9.5, subkategori abelian dan teorema rekonstruksi.

% segment-id: o014.li2.chapter9.reconstruction-b.p001
Hubungan hasil di atas dengan subkategori abelian juga mudah dijelaskan.

% segment-id: o014.li2.chapter9.reconstruction-b.l001
\begin{lemma}\label{prop:subcat-gen-L}
	Misalkan $\mathcal{A}$ dan $\omega$ memenuhi syarat-syarat
	Lema~\ref{prop:Tannaka-projective-generator}. Setiap subkategori abelian
	$\mathcal{A}'$, bersama $\omega|_{\mathcal{A}'}$, juga memenuhi
	syarat-syarat tersebut. Selain itu, homomorfisme kanonik
	$\coEnd(\omega|_{\mathcal{A}'})\to\coEnd(\omega)$ dari
	Catatan~\ref{rem:L-cogebra-functoriality} bersifat monik.
\end{lemma}
\begin{proof}
	Jelas bahwa $\mathcal{A}'$ hingga secara lokal. Tanpa mengurangi keumuman, kita
	boleh mengandaikan bahwa $\Obj(\mathcal{A}')\subset\Obj(\mathcal{A})$
	tertutup terhadap isomorfisme. Kita memakai sifat-sifat umum berikut;
	argumennya tidak sulit dan dapat ditemukan dalam latihan
	\CHref{sec:app-Abel}.
	\begin{itemize}
		\item Pilih generator projektif $s$ bagi $\mathcal{A}$. Objek ini
		memiliki objek hasil bagi terbesar yang terletak dalam
		$\mathcal{A}'$, dinotasikan oleh $s'$, dan $s'$ merupakan generator
		projektif bagi $\mathcal{A}'$.
		\item Tuliskan $A:=\End(s)$ dan
		$A':=\End(s')\simeq\Hom(s,s')$. Terdapat ideal dua sisi
		$\mathfrak a$ sedemikian rupa sehingga $A\to A'$ menginduksi
		$A/\mathfrak a\rightiso A'$. Lebih eksplisit, jika
		$t:=\Ker[s\twoheadrightarrow s']\subset s$, maka
		$\mathfrak a=\Hom(s,t)\subset A$.
		\item Setelah $\mathcal{A}$ diidentifikasikan dengan
		$\catesubd{Mod}{\mathrm{fg}}A$, $\mathcal{A}'$ teridentifikasi dengan
		subkategori penuh yang dibentuk oleh modul-modul yang dianihilasi oleh
		$\mathfrak a$.
	\end{itemize}

	Pilih basis $\phi_1,\ldots,\phi_m$ dari ruang vektor-$\Bbbk$
	$\Hom(s,t)$. Sifat generator memberikan barisan eksak
	% segment-id: o014.li2.chapter9.reconstruction-b.q001
	\[ s^{\oplus m} \xrightarrow{(\phi_i)_i} s \to s' \to 0. \]
	Tuliskan $P:=\omega(s)$ dan $P':=\omega(s')$. Aljabar $A$ bekerja dari
	kiri pada $P$ melalui $\omega$, dan barisan eksak yang bersesuaian
	% segment-id: o014.li2.chapter9.reconstruction-b.q002
	\[ P^{\oplus m} \xrightarrow{(\omega(\phi_i))_i} P \to P' \to 0 \]
	menunjukkan bahwa $P'\simeq P/\mathfrak aP$.

	Dalam bukti Lema~\ref{prop:Tannaka-projective-generator} telah terlihat
	bahwa $(\mathord\cdot)\dotimes{A}P:\cated{Mod}A\to\cated{Mod}B$ eksak.
	Bersama deskripsi di atas, ini memberikan
	% segment-id: o014.li2.chapter9.reconstruction-b.q003
	\begin{align*}
		(P')^\vee \dotimes{A/\mathfrak{a}} P' & \simeq (A/\mathfrak{a} \dotimes{A} P)^\vee \dotimes{A/\mathfrak{a}} ( A/\mathfrak{a} \dotimes{A} P) \\
		& \simeq (A/\mathfrak{a} \dotimes{A} P)^\vee \dotimes{A} P \\
		& \hookrightarrow P^\vee \dotimes{A} P.
	\end{align*}
	Tidak sulit untuk memeriksa bahwa ini tepat homomorfisme kanonik
	$\coEnd(\omega|_{\mathcal{A}'})\to\coEnd(\omega)$.
\end{proof}

% segment-id: o014.li2.chapter9.reconstruction-b.eg001
\begin{example}\label{eg:Tannakian-I-graded-0}
	Berikut beberapa penerapan awal
	Lema~\ref{prop:Tannaka-projective-generator} untuk $B=\Bbbk$; ingat
	bahwa $\Bbbk$ diasumsikan medan. Untuk himpunan tak kosong $I$, definisikan
	$\cate{Vect}_I(\Bbbk)$ sebagai kategori ruang vektor-$\Bbbk$ bergradasi
	$I$ (Contoh~\ref{eg:graded-module}, dengan $\epsilon$ trivial). Objek
	$(M^i)_{i\in I}$ yang memenuhi
	$\sum_i\dim_{\Bbbk}M^i<\infty$ membentuk subkategori penuh
	$\mathcal{A}:=\cate{Vect}_{I,\mathrm f}(\Bbbk)$. Tuliskan $\Bbbk[I]$
	untuk ruang vektor-$\Bbbk$ berbasis $I$; struktur berikut menjadikannya
	koaljabar kokomutatif:
	% segment-id: o014.li2.chapter9.reconstruction-b.q004
	\[ \epsilon(i) = 1, \quad \Delta(i) = i \otimes i, \quad i \in I. \]
	Definisikan funktor
	$\omega_I:\mathcal{A}\to\cate{Vect}_{\mathrm f}(\Bbbk)$ dengan
	memetakan $(M^i)_{i\in I}$ ke $\bigoplus_iM^i$.
	\begin{enumerate}
		\item Kasus paling sederhana $|I|=1$ sama dengan mengambil
		$\mathcal{A}=\cate{Vect}_{\mathrm f}(\Bbbk)$ dan
		$\omega_I=\identity$; generator projektif $s$ dapat diambil sebagai
		$\Bbbk$. Perhitungan langsung dari definisi, atau
		Lema~\ref{prop:Tannaka-projective-generator}(i), memberikan
		% segment-id: o014.li2.chapter9.reconstruction-b.q005
		\[ \coEnd(\omega_I) \simeq \Bbbk, \quad \epsilon(1) = 1, \quad \Delta(1) = 1 \otimes 1. \]
		Dengan demikian, bila $|I|=1$, berlaku
		$\coEnd(\omega_I)\simeq\Bbbk\simeq\Bbbk[I]$.

		\item Berikutnya, jika $I$ hingga, maka
		$\coEnd(\omega_I)\simeq\Bbbk[I]$. Ini dapat diturunkan dari kasus
		sebelumnya atau dipahami dengan menerapkan
		Lema~\ref{prop:Tannaka-projective-generator}(i), dengan mengambil
		generator projektif $s$ yang ditentukan oleh $M^i=\Bbbk$ untuk semua
		$i$.

		Perhatikan bahwa jika $I\subset J$, pembenaman
		$\cate{Vect}_{I,\mathrm f}(\Bbbk)\subset
		\cate{Vect}_{J,\mathrm f}(\Bbbk)$ menginduksi
		$\Bbbk[I]\to\Bbbk[J]$, yakni pembenaman koaljabar yang diinduksi oleh
		$I\hookrightarrow J$.

		\item Untuk setiap himpunan tak kosong $I$, berlaku pula
		$\coEnd(\omega_I)\simeq\Bbbk[I]$. Lebih tepatnya, jika
		$x\in M^i$ dan $\phi\in(M^i)^\vee$ memenuhi $\phi(x)=1$, maka
		$[\phi\otimes x]\in\coEnd(\omega_I)$ bersesuaian dengan
		$i\in\Bbbk[I]$.

		Himpunan $I$ merupakan gabungan terfilter dari subhimpunan hingga,
		dan kategori $\cate{Vect}_{I,\mathrm f}(\Bbbk)$ serta koaljabar
		$\Bbbk[I]$ juga direduksi ke kasus hingga dengan cara yang sama.
		Karena kasus $I$ hingga telah diketahui,
		Proposisi~\ref{prop:coHom-colim} langsung memberikan
		% segment-id: o014.li2.chapter9.reconstruction-b.q006
		\[ \coEnd(\omega_I) \simeq \varinjlim_{J \subset I, |J| < \infty} \coEnd(\omega_J). \]
		Pengambilan kolimit adalah teknik kunci yang akan muncul kembali dalam
		bukti berikutnya.
	\end{enumerate}
\end{example}

% segment-id: o014.li2.chapter9.reconstruction-b.p002
Teorema berikut merupakan hasil antara utama sejauh ini, dan buktinya
bergantung pada hasil teknis dari \S\ref{sec:locally-finite-Abel-cat}.
Ingat bahwa jika $(L,\mu,\eta,\Delta,\epsilon)$ bialjabar-$B$,
Proposisi~\ref{prop:bialgebra-Mod-monoidal} membekali
$\cated{Comod}L$ dengan struktur monoidal alami: operasi $\otimes$ berasal
dari $\mu$, dan unitnya ialah $B$ dengan koaksi yang diberikan oleh
$\eta:B\to B\dotimes{B}L\simeq L$. Perhatikan bahwa
$\catesubd{Comod}{\mathrm{pf}}L$ merupakan subkategori monoidalnya.

% segment-id: o014.li2.chapter9.reconstruction-b.th001
\begin{theorem}\label{prop:Tannaka-aux1}
	\index{teorema rekonstruksi@teorema rekonstruksi (Reconstruction Theorem)}
	Misalkan $\Bbbk$ medan, $B$ aljabar-$\Bbbk$, $\mathcal{A}$ kategori
	abelian hingga secara lokal, dan
	$\omega:\mathcal{A}\to\cated{Mod}B$ funktor setia eksak yang mengambil
	nilai dalam $\catesubd{Mod}{\mathrm{pfg}}B$.
	\begin{enumerate}[(i)]
		\item Funktor $\overline\omega$ dalam faktorisasi kanonik
		\eqref{eqn:L-canonical-decomp} merupakan ekuivalensi kategori. Selain
		itu, terdapat keluarga subkategori abelian $\mathcal{A}'$ dari
		$\mathcal{A}$ sedemikian rupa sehingga
		\begin{compactitem}
			\item $\bigcup_{\mathcal{A}'}\mathcal{A}'=\mathcal{A}$ merupakan
			gabungan terfilter subkategori;
			\item setiap $\mathcal{A}'$ memiliki generator projektif,
		\end{compactitem}
		dan untuk keluarga tersebut terdapat isomorfisme koaljabar
		% segment-id: o014.li2.chapter9.reconstruction-b.q007
		\[ \coEnd(\omega) \simeq \varinjlim_{\mathcal{A}'} \coEnd(\omega'), \quad \omega' := \omega|_{\mathcal{A}'}. \]
		Homomorfisme transisi di ruas kanan berasal dari
		Catatan~\ref{rem:L-cogebra-functoriality} dan semuanya merupakan pembenaman
		koaljabar. Dalam situasi ini,
		% segment-id: o014.li2.chapter9.reconstruction-b.q008
		\[ \catesubd{Comod}{\mathrm f}\coEnd(\omega)=
		\catesubd{Comod}{\mathrm{pf}}\coEnd(\omega). \]

		\item Jika $\mathcal{A}$ memiliki struktur kategori monoidal,
		$B=\Bbbk$, dan $\omega$ funktor monoidal, maka terhadap struktur
		bialjabar-$\Bbbk$ pada $\coEnd(\omega)$ dari
		Proposisi~\ref{prop:L-bialgebra}, kedua funktor
		$\overline\omega$ dan $U$ bersifat monoidal.
		\item Dengan asumsi yang sama, jika lebih lanjut $\mathcal{A}$ kategori
		monoidal simetris dan $\omega$ kompatibel dengan struktur kepang, maka
		$\coEnd(\omega)$ komutatif sebagai aljabar, sedangkan
		$\overline\omega$ dan $U$ juga kompatibel dengan struktur kepang.
	\end{enumerate}
\end{theorem}
\begin{proof}
	Proposisi~\ref{prop:subcat-union} menuliskan $\mathcal{A}$ sebagai
	gabungan terfilter subkategori abelian $\mathcal{A}'$ yang memiliki
	generator projektif; misalnya, kita dapat mengambil
	$\mathcal{A}'=\lrangle{X}$ untuk $X\in\Obj(\mathcal{A})$.
	Proposisi~\ref{prop:coHom-colim} menunjukkan bahwa
	$\coEnd(\omega)\simeq\varinjlim_{\mathcal{A}'}\coEnd(\omega')$, dan
	Lema~\ref{prop:subcat-gen-L} menunjukkan bahwa semua homomorfisme
	transisi adalah pembenaman koaljabar. Faktorisasi kanonik
	\eqref{eqn:L-canonical-decomp} diperoleh dengan mengambil gabungan
	terfilter dari
	% segment-id: o014.li2.chapter9.reconstruction-b.q009
	\[\begin{tikzcd}
		\mathcal{A}' \arrow[r, "{\overline{\omega'}}"] & \catesubd{Comod}{\mathrm{pf}}\coEnd(\omega') \arrow[r, "U"] & \catesubd{Mod}{\mathrm{pfg}}B.
	\end{tikzcd}\]
	Lema~\ref{prop:Tannaka-projective-generator} menyatakan bahwa setiap
	$\overline{\omega'}$ merupakan ekuivalensi, sehingga $\overline\omega$ juga
	merupakan ekuivalensi. Selain itu,
	$\catesubd{Comod}{\mathrm{pf}}\coEnd(\omega')=
	\catesubd{Comod}{\mathrm f}\coEnd(\omega')$; proses yang sama sekaligus
	memberikan kesetaraan tersebut untuk $\coEnd(\omega)$. Ini membuktikan
	(i).

	Jika $\mathcal{A}$ monoidal dan $B=\Bbbk$, funktor lupa $U$ selalu
	monoidal terhadap struktur pada $\cated{Comod}\coEnd(\omega)$ yang
	berasal dari bialjabar $\coEnd(\omega)$. Diagram komutatif
	\eqref{eqn:L-bialgebra-prod}, dengan semua
	$\omega_i,\omega_i',\omega_i''$ diganti oleh $\omega$, menunjukkan bahwa
	$\overline\omega$ juga monoidal. Ini membuktikan (ii), dan argumen serupa
	membuktikan (iii).
\end{proof}

% segment-id: o014.li2.chapter9.reconstruction-b.eg002
\begin{example}\label{eg:Tannakian-I-graded}
	Misalkan $\Gamma$ monoid dan $B=\Bbbk$ sebuah medan, sehingga
	$\catesubd{Mod}{\mathrm{pfg}}B=\cate{Vect}_{\mathrm f}(\Bbbk)$.
	Contoh~\ref{eg:Tannakian-I-graded-0} menunjukkan bahwa koaljabar yang
	bersesuaian dengan
	% segment-id: o014.li2.chapter9.reconstruction-b.q010
	\[ \omega_\Gamma: \cate{Vect}_{\Gamma, \mathrm{f}}(\Bbbk) \to \cate{Vect}_{\mathrm{f}}(\Bbbk), \quad (M^\gamma)_{\gamma \in \Gamma} \to \bigoplus_\gamma M^\gamma \]
	ialah $\Bbbk[\Gamma]$. Selain itu, $\Bbbk[\Gamma]$ membawa
	struktur bialjabar (Contoh~\ref{eg:monoid-bialgebra}). Di sisi lain,
	$\cate{Vect}_{\Gamma,\mathrm f}(\Bbbk)$ kategori monoidal dan
	$\omega_\Gamma$ funktor monoidal (Contoh~\ref{eg:graded-module}); ini
	juga membekali $\Bbbk[\Gamma]$ dengan struktur bialjabar. Kedua struktur
	tersebut sama, sebagaimana dapat diperiksa dengan membandingkan deskripsi
	perkalian dalam \eqref{eqn:L-tensor} dengan
	Contoh~\ref{eg:Tannakian-I-graded-0}; perinciannya diserahkan kepada
	pembaca.
\end{example}
